\documentclass[11pt,letter]{article}

\usepackage[margin=2.5cm]{geometry}
\usepackage{amsmath}  % \eqref{eq:ratio}
\usepackage{amssymb}
\usepackage{bm}
\usepackage{booktabs}
\usepackage{array}
\usepackage{ragged2e}
\usepackage{seqsplit}
\usepackage{enumitem}
\usepackage[hidelinks]{hyperref}
\usepackage{cleveref}  % \cref or \Cref{sec:ratio}   Add after hyperref.
\usepackage{setspace}
\setstretch{1.15}


\setlength{\emergencystretch}{3em}
\newcolumntype{L}[1]{>{\RaggedRight\arraybackslash}p{#1}}

%----------------------------------------------------------------------
%  Text helpers
%----------------------------------------------------------------------
\newcommand{\file}[1]{\texttt{\seqsplit{#1}}}
\newcommand{\ToClaude}[1]{[\textbf{CLAUDE}: {#1}]}
\newcommand{\COMMENT}[1]{}

%----------------------------------------------------------------------
%  Notation.  Lowercase italic = what is physically there (never
%  computed).  Bold uppercase = an image computed from a photo.  Plain =
%  one number.  Prime = the pre-image.
%----------------------------------------------------------------------
\newcommand{\iE}{\mathbf{E}}   % reading as stored, DN
\newcommand{\iD}{\mathbf{D}}   % light counts, DN
\newcommand{\iP}{\mathbf{P}}   % shape of the stage's light
\newcommand{\iG}{\mathbf{G}}   % gray level map, DN
\newcommand{\iR}{\mathbf{R}}   % reflectance relative to the gray chips
\newcommand{\iW}{\mathbf{W}}   % pre-image white, lined up
\newcommand{\tW}{\widetilde{\mathbf{W}}}   % W smoothed
\newcommand{\hW}{\widehat{\mathbf{W}}}     % estimated white
\newcommand{\iT}{\mathbf{T}}   % fraction the dye let through
\newcommand{\iA}{\mathbf{A}}   % stain absorbance (the output)
\newcommand{\tblA}{\mathbb{A}} % A for all channels
\newcommand{\vgamma}{\bm{\gamma}}
\newcommand{\reg}[2]{\underset{\text{#1}}{#2}}
\newcommand{\avg}[2]{\overline{\reg{#1}{#2}}}
\newcommand{\MD}{\mathcal{D}}
\newcommand{\Lten}{\log_{10}}


\title{From photons to stain absorbance\\
\large The chain for one pixel, derived from one thought experiment}
\author{Colorimeter project --- pad analysis}
\date{Draft for Bill's review, 2026-10-05.  Companion to
\file{analysis\_chain.tex}; no code is changed by this document.}

\begin{document}
\maketitle

\begin{abstract}
\noindent
A photo of the pad records light that the lamp sent out, that reached the
pad, and that the pad sent back.  Writing that down for one pixel gives a
product of factors: the lamp, the shape of the light on the stage, the
pad's shadow, the paper, and the dye.  This document follows that product
through each step of the chain and shows, one line of algebra at a time,
which factor each step removes, until only the dye is left.  It then shows
why the pre-image should be both used as it is and smoothed (method
$\MD$, with both weights non-zero), and why the change from reflectance to
absorbance belongs at the very end.
\end{abstract}

\tableofcontents

%======================================================================
\section{What this document does}
\label{sec:scope}

It follows one pixel of one channel (one LED and one Bayer layer) from the
stored photo to the stain's absorbance.  The two photos are simplified to
what matters for the algebra:
\begin{itemize}[nosep]
  \item the \textbf{pre-image} (the strip before exposure to arsine) shows
    the white pad and its shadow, with no dye;
  \item the \textbf{post-image} (after exposure) shows the same pad, with
    the same shadow, and dye.
\end{itemize}
It leaves out everything that does not change the algebra: choosing the
exposure time, hot pixels, finding the pad, how the lining-up shift is
found, how the smoothing is done near the pad's edge, the yellow/brown
split, and the calibration against arsenic.

\medskip
\noindent The chain uses three ways of expressing the light, each for a
reason:

\medskip
\begin{tabular}{@{}L{2.6cm}L{2.0cm}L{10.0cm}@{}}
\toprule
Domain & Steps & What it is used for, and why \\
\midrule
counts (DN) & 1 & Removing the dark reading, which \emph{adds} to the
  light, so it must be subtracted from the raw counts. \\
reflectance & 2--5 & Removing factors (lamp, stage, shadow, paper) by
  division, and averaging light (lining up, smoothing), which must be done
  on light itself (\Cref{sec:rules}). \\
absorbance & 6 & The output.  It grows in proportion to the amount of dye,
  and two dyes' absorbances add, which is what the yellow/brown split
  needs. \\
\bottomrule
\end{tabular}

%======================================================================
\section{Simplified outline of the signal path}
\label{sec:thought}

The measured signal is a count of photons landing on the camera sensors.  
From that we extract an estimate of the dye on the paper indicator pad, 
and the dye is later related to a chemical concentration. 
Here we follow the path of the photons from the LEDs to a camera sensor and show how the dye amount is measured at one pixel of the camera image.
\begin{itemize}[nosep]
  \item The LEDs emit photons toward the stage. %, on which are a paper indicator pad and a reference gray chip.
    The camera focuses a sensor (an individual pixel) on a small spot the stage, on either the paper indicator pad or a reference gray chip. 
    A fraction of the photons that hit that spot are reflected back to a sensor, where they are counted and reported as a pixel value.
    The maximum number of photons, $n$, that can be reflected back to the sensor is set by the staging, the lamp's brightness, and the exposure time.  
    The staging is fixed and is subsumed in the concept of brightness.
  \item The light hitting the spot can be obstructed by shading.
    When one or more LEDs cannot see the spot because of shadows, then only a fraction $b$ of the light is added, and the remaining fraction of the light is not observed.
   	$b = 1$ when the spot is fully lit, and $b < 1$ when the spot is partly shaded.
   	The number of photons that can be reflected from the spot is $b \cdot n$.
  \item $a$ is the surface reflectance, the fraction of photons reflected back from the spot to the sensor.  
    On the pad, $a$ is a constant property of the paper, plus the effect of any dye on the paper that absorbs photons. 
    A shadow does not change $a$.
    The reflected photons that are captured by a sensor are then $d = a \cdot b \cdot n$
  \item $e$ is a baseline signal that the sensor always reports. 
    It is called the \textit{dark level} because it is the value of $d$ when $n=0$.
    The sensor measurment is the sum of the counted photons plus the dark level.
    \begin{equation}
    	d_\text{raw, lit} = a \cdot b \cdot n + e
    \end{equation}
    %$d_\text{raw, lit} = a \cdot b \cdot n + e$. 
\end{itemize}

For each measurement 3 photos are taken:  
a dark image $e$, 
a lit image taken before the pad is exposed to chemical, $d'$, 
and a final image taken after chemical exposure, $d$. 
The latter two are called the \textit{pre-image} and the \textit{post-image}. 
We write a prime ($'$) to distinguish the pre-image. 

\subsection{Dark level removal}

The dark level $e$ is removed by measuring it in an unlit photo, and subtracting it from the raw sensor reading in a lit photo.
\begin{alignat}{2}
	e & = d_\text{raw, unlit}  \\
	d_\text{lit} & = d_\text{raw, lit} - e
\end{alignat}

The dark level varies with the temperature of the camera, so it is remeasured each time that the camera is used.
In practice, the change in $e$ over a short sequence of photos is negligible.

The dark level varies by pixel, but the variation is mostly noise. 
The best predictor of dark level for a given pixel in a given image is the mean dark level of all valid pixels in a contemporaneous dark image.   

A small percentage of images report dark levels that are identified as outliers.
We set a threshold range and label pixels outside of the threshold range as invalid.

\subsection{Shading removal}

Shade is removed from a post-image of the pad, $d$, by comparing it with a pre-image, $d'$ of the same pad.
%The same spot is measured in two photos, a pre-image taken before dye is developed, and a post-image taken after dye is developed.  
%Write a prime ($'$) to denote the pre-image.  
Assume that the shadow $b$ is unchanged, $b' = b$; 
But the surface changes from $a'$ to $a$; 
and the exposure changes from $n'$ to $n$.
We have
\begin{equation}
  d' = a'\, b\, n', \qquad d = a\, b\, n .
  \label{eq:te}
\end{equation}

\noindent\textbf{The ratio} of $d / d'$ cancels the shading, independent of the exposure and the amount of dye 
\begin{equation}
  \frac{d}{d'} = \frac{a\, b\, n}{a'\, b\, n'}
  = \frac{a}{a'} \cdot \frac{n}{n'} .
  \label{eq:te-ratio}
\end{equation}
Shading $b$ is gone;  
leaving the change in the dye, $a/a'$, 
and the change in the exposure, $n/n'$.

We later remove two kinds of shading,  
one due to paper roughness and the other due to fiber disturbance along the pad edges.
Their effect and removal are the same, but they are measured differently.

\subsection{Balancing the exposure}

Equation \eqref{eq:te-ratio} contains unknowns $n$ and $n'$ that we wish to remove.
They are unknown because the exposure time and LED emissions may both vary.
We can remove the ratio $n/n'$ by observing reflectance from a reference gray chip with invariant reflectance and shading.
The fraction of light returned by the gray chip in every image is $\rho$, which replaces $a \cdot b$.

Measure the dark-corrected reflectance on each pre-image and each post-image, $d_g'$ and $d_g$.
There is no shading on the chip, and there is no change in the chip color between images, 
so the measures are only a function of $\rho \, n$, or $\rho' n'$
\begin{equation}
	d_g' = \rho\, n_g', \qquad d_g = \rho \, n_g .
\end{equation}
 
Taking the ratio $d_g / d_g'$ gives
\begin{equation}
	\frac{d_g}{d_g'} = \frac{\rho \, n}{\rho \, n'} = \frac{n}{n'}
	\label{eq:dggdash}
\end{equation}
We can substitute \eqref{eq:dggdash} for  $n / n'$ in \eqref{eq:te-ratio}
and solve for $a / a'$
\begin{equation}
	\frac{a}{a'} = \frac{d}{d'} \frac{d_g'}{d_g}
	\label{eq:daadash}
\end{equation}

\subsection{Finding $d_g$ on the pad}

\COMMENT{In this section I need to show clearly and simply how $P$ works to estimate $d_g'$ and $d_g$ during measurments.  I think I have to show the basic solution of $P$ during calibration, minimizing $d_g - P$ using least squares, and then how $P$ is used to extrapolate from the reference gray chips to the equivalent reference gray value at the pad during measurements.  I'd like to first show how that would work using a single gray reference chip with only a scalar multiplier on $P$, and then how it would work using multiple gray reference readings and linear or bilinear multiplier on $P$.  The difficulty is to show that the algebra is correct without getting too detailed. For example, I'd like to say that $P$ is found by minimizing $(P - d_g)^2$ over the area, and not show the matrix solution.  }

The lighting changes across the stage, and the values $d_g$ measured on a reference gray-chip near the pad will differ from what would be read from a reference gray-chip at the pad.
In a calibration phase, a matching gray chip is placed in the pad location, and a polynomial $P(i,j)$ is found 
that minimizes $[P(i,j) - d_{g,(i,j)}]^2$ over all pixels $(i,j)$ over the area of the pad and the reference gray chip.
It is also determined during calibration that across various exposure times the value of $P$ varies only by a constant, $\varphi$. 
We therefore have for pre-image and post-image during measurments, 
\begin{equation} \label{eq:dgofPprepost1}
	d'_{g,(i,j)} = \varphi' P_{i,j} \qquad \ d_{g,(i,j)} = \varphi P_{i,j} 
\end{equation}
Taking ratios of the left and right hand sides of equations \eqref{eq:dgofPprepost1} gives
\begin{equation} \label{eq:dgofPprepost2}
	\frac{d_{g,(i,j)}}{d'_{g,(i,j)}} = \frac{\varphi P_{i,j}}{\varphi' P_{i,j}} 
\end{equation}
We se that $P_{(i,j)}$ cancels, and that the ratio $d_g/d_g'$ is equal to a constant $\Phi = \varphi/\varphi'$.
In a pre-image and post-image pair, we can estimate the constant $\Phi$ by sampling $d_g$ and $d_g'$ at the reference gray chip
\begin{equation} \label{eq:dgofPprepost3}
	\widehat\Phi = \text{mean}\left[\frac{d_{g,(i,j)}}{d'_{g,(i,j)}}\right] \quad (i,j)\in\text{ gray chip}
\end{equation}
Then for any pixel $(i,j)$ on the pad we can write \eqref{eq:daadash} as 
\begin{equation}
	\frac{a}{a'} = \frac{d}{d'} \widehat\Phi
	\label{eq:daadashPhi}
\end{equation}

It remains to show that $\Phi$ really is a constant across the stage, 
and that $P$ really is invariant across images.

\COMMENT{
that models $d_g$ across the areas of the pad and the reference chip. 
For a pixel at $(i,j)$, we have
\begin{equation} \label{eq:dgofPvarphi}
	d_{g,(i,j)} = \varphi P_{i,j} % n_{(i,j)} 
\end{equation}
where $\varphi$ is some varying factor due to differences in LED brightness and exposure duration.
We anticipate that $P_{i,j}$ should be fixed
because the geometry of the LEDs and the reference gray chip is held invariant after calibration.
This is verified by showing $P$ can be fit satisfactorily  during calibration using a scalar $\varphi$ over different exposure durations.
%
With $P$ known, we write for the pre-image and post-image
\begin{equation} \label{eq:dgofPprepost}
	d'_{g,(i,j)} = \varphi' P_{i,j} \qquad \ d_{g,(i,j)} = \varphi P_{i,j} 
\end{equation}
So that 
%
that \eqref{eq:dgofPvarphi} provides a good estimate of $d_{g,(i,j)}$ using a fixed $P$ and varying $\varphi$ over many photos with varying exposure times.
%
 that the $d_{g,(i,j)}$ 
and we assume that $P$ measured during calibration and $P$ measured during 
%
 matches the average relative reflectance from the reference gray chips in their fixed locations on the stage, and on the gray chip in the pad location.
The  gray chip level can be estimated at any point on the pad by fitting the polynomial to match the observed reference gray chip levels.
}

%%%%%%%%%%%%%%%%%%%%%%%%%%%%%%%%%%%%%%%%%%%%%%%%%%%%%%%%%%%%%%%%%%%%%%%%
%%%%%%%%%%%%%%%%%%%%%%%%%%%%%%%%%%%%%%%%%%%%%%%%%%%%%%%%%%%%%%%%%%%%%%%%
\section{END OF chain\_derivation2.tex}
Beyond here is not edited.

\end{document}

Section dye effect is slightly edited
\section{dye effect} 

We assume that light absorbance by the dye behaves approximately as light transmitted through a dyed liquid.
In a dyed liquid, the dye acts as a filter that transmits a fraction $t$ of the light.  
On the paper pad, the dyed surface reflects light, and we model the dye as absorbing some, permitting the fraction $t$ to be reflected.  
On the undyed paper $t$ = 1.
On the dyed paper, the reflected amount is 
\begin{equation}
	a  = a' \, t
	\label{eq:dye1}
\end{equation}
so the final measure of the dye
\begin{equation}
	t   = \frac{a}{a'}
	\label{eq:dye2}
\end{equation}

Combining \eqref{eq:daadashPhi} and \eqref{eq:dye2}, gives $t$ in terms of quantities measured in the pre-image and post-image pair
\begin{equation}
	t   = \frac{d}{d'} \widehat\Phi
	\label{eq:dye3}
\end{equation}


\section{dye amount and arsenic concentration}

The Beer-Lambert law for light transmission through a dyed liquid states that the 

\end{document}
%%%%%%%%%%%%%%%%%%%%%%%%%%%%%%%%%%%%%%%%%%%%%%%%%%%%%%%%%%%%%%%%%%%%%%%%
%%%%%%%%%%%%%%%%%%%%%%%%%%%%%%%%%%%%%%%%%%%%%%%%%%%%%%%%%%%%%%%%%%%%%%%%
The exposures $n'$ and $n$,  may differ
, which the gray chips
will measure (\Cref{sec:ratio}).
%
\noindent\textbf{The difference} does not cancel the shadow:
\begin{equation}
  d - d' = b\,\bigl(a\, n - a'\, n'\bigr).
  \label{eq:te-diff}
\end{equation}
$b$ stays as a multiplier, so the shadow stays in the answer.
%
\noindent\textbf{Example.}  Equal exposures ($n = n' = 1$); the stain
halves the light ($a' = 1$, $a = 0.5$); in the shadow a quarter of the
light is blocked ($b = 0.75$):
%
\medskip
\begin{center}
\begin{tabular}{@{}lrrrr@{}}
\toprule
 & $d'$ (pre) & $d$ (post) & ratio $d/d'$ & difference $d - d'$ \\
\midrule
fully lit ($b = 1$)    & 1.00 & 0.50  & \textbf{0.50} & $-0.50$ \\
in shadow ($b = 0.75$) & 0.75 & 0.375 & \textbf{0.50} & $-0.375$ \\
\bottomrule
\end{tabular}
\end{center}
%
\noindent The ratio reads the same stain in and out of the shadow.  The
difference reads the shadowed spot as 25\,\% less stained.
%
\subsection{Two rules}
\label{sec:rules}

\textbf{Rule 1: combining light is an ordinary (arithmetic) average or sum,
of counts or of reflectance.}  Light adds: from different LEDs (that is
what $b$ is), from different places, and from different moments.  A pixel
records the light collected over its area.  So any step that stands for
``the light over a shifted or a larger area'' --- lining up, smoothing ---
must be an arithmetic average of light.

\textbf{Rule 2: an influence that enters as a factor is removed by
division, which is exactly a subtraction in absorbance.}  The logarithm of
a product is the sum of the logarithms:
\begin{equation}
  -\Lten d = -\Lten a - \Lten b - \Lten n .
  \label{eq:log}
\end{equation}
Dividing $d$ by $b$, or subtracting $-\Lten b$ from $-\Lten d$, gives the
same number.  So for a division, it does not matter which domain it is
written in.  For an average, it does: by rule 1, averages are done on light.

%======================================================================
\section{Notation}
\label{sec:notation}

\begin{tabular}{@{}L{3.4cm}L{11.4cm}@{}}
\toprule
Style & Meaning \\
\midrule
$a$, $b_{ij}$ (lowercase italic) & What is physically there: the
  thought experiment's factors.  They are never computed from a photo, so
  they have no bold form. \\
$\iR$ (bold uppercase) & An image computed from a photo: one value per
  pixel, one channel. \\
$R_{ij}$ (plain, subscripted) & One element of that image: row $i$,
  column $j$.  A single number. \\
$\vgamma$ (bold lowercase) & A list of numbers, here one per column. \\
$\tblA$ (blackboard bold) & A table: one row per pixel, one column per
  channel. \\
$e$, $\lambda$ (plain) & One number. \\
$'$ (prime) & The pre-image.  Same letter as in the post-image. \\
$\avg{g1}{D}$ & The mean over a region; the region is named underneath
  (g1, g2 = the left and right gray chips). \\
$\widetilde W$, $\widehat W$ & Tilde: smoothed.  Hat: estimated. \\
\bottomrule
\end{tabular}

\medskip
\noindent Arithmetic between images is done pixel by pixel; no matrix
product is meant anywhere.  Equations are written for one element.

\subsection*{What is physically there}

\begin{tabular}{@{}L{2.4cm}L{12.4cm}@{}}
\toprule
$n_j$, $n'_j$ & The light level of the photo at column $j$: exposure time
  $\times$ lamp brightness $\times$ the camera's counts per photon.  It may
  change slowly from left to right. \\
$P_{ij}$ & The shape of the light on the stage (the flat-field polynomial),
  mean 1.  Measured once; the same for both photos. \\
$b_{ij}$ & The shadow: the fraction of the stage's light that reaches this
  spot of the pad.  It moves with the strip. \\
$a_{ij}$, $a'_{ij}$ & The fraction of the arriving light that the surface
  sends back, after and before exposure. \\
$t_{ij}$ & The fraction of the light that the dye lets through; 1 where
  there is no dye. \\
$\alpha_{ij}$ & The amount of dye on the absorbance scale,
  $\alpha_{ij} = -\Lten t_{ij}$.  Bigger means more dye. \\
$\rho$ & The gray chips' own return fraction.  The chips are fixed to the
  stage, so $\rho$ is the same in both photos. \\
\bottomrule
\end{tabular}

\subsection*{What is computed from the photos}

\begin{tabular}{@{}L{3.4cm}L{9.6cm}L{1.6cm}@{}}
\toprule
Symbol & Meaning & Step \\
\midrule
$\iE$, $\iE'$ & the reading as stored, DN & 1 \\
$e$, $e'$ & the dark level: one number per Bayer layer per photo & 1 \\
$\iD$, $\iD'$ & light counts, $\iE - e$; the image form of $d$ & 1 \\
$\gamma_1$, $\gamma_2$ & the light level measured at each gray chip & 2 \\
$f_j$, $\gamma_j$ & fraction of the way from the left chip to the right;
  the level blended between the chips & 2 \\
$\iG$, $\iG'$ & the gray level at each pixel & 2 \\
$\iR$, $\iR'$ & reflectance relative to the gray chips, each photo in its
  own frame & 2 \\
$\iW$, $\eta_q$ & the pre-image's $\iR'$ lined up on the post-image's
  grid: the white, one value per pixel; $\eta_q$ are the lining-up weights
  & 3 \\
$\iT$ & the fraction the dye let through, as measured & 4, 5 \\
$\tW$, $k_m$ & $\iW$ smoothed; $k_m$ are the smoothing weights & 5 \\
$\hW$, $\lambda$ & the estimated white, a blend of $\iW$ and $\tW$;
  $\lambda$ is the weight on $\iW$ & 5 \\
$\delta$, $\varepsilon$, $\tau$ & fractional errors: post-image noise,
  pre-image noise, pattern lost by smoothing & 5 \\
$\iA$ & \textbf{the stain's absorbance: the output} & 6 \\
$\tblA$ & $\iA$ for all channels: pixels $\times$ channels & 6 \\
\bottomrule
\end{tabular}

%======================================================================
\section{A pixel's reading}
\label{sec:model}

The thought experiment, written for pixel $(i, j)$ of a real photo, needs
two more factors and one added term.  The light reaching the stage is not
the same everywhere: it is $n_j P_{ij}$, a level $n_j$ for this photo times
a fixed shape $P_{ij}$.  And the camera reads a little even with no light,
the dark level $e$.  So the reading at a pad pixel in the post-image is
\begin{equation}
  E_{ij} = a_{ij}\, b_{ij}\, n_j\, P_{ij} + e ,
  \label{eq:model-pad}
\end{equation}
at a gray-chip pixel, where the surface is the chip and nothing shades it,
\begin{equation}
  E_{ij} = \rho\, n_j\, P_{ij} + e ,
  \label{eq:model-chip}
\end{equation}
and in the pre-image, in its own frame,
\begin{equation}
  E'_{ij} = a'_{ij}\, b_{ij}\, n'_j\, P_{ij} + e' , \qquad
  E'_{ij} = \rho\, n'_j\, P_{ij} + e' \ \text{(chip)}.
  \label{eq:model-pre}
\end{equation}
The dye acts as a filter over the paper: the post-image surface returns
what the paper returned, times the fraction the dye lets through,
\begin{equation}
  a_{ij} = a'_{ij}\, t_{ij} .
  \label{eq:dye}
\end{equation}

\subsection*{Assumptions}

Each step below names the assumption it uses.
\begin{enumerate}[label=(A\arabic*), leftmargin=3.2em]
  \item \textbf{The dark reading adds} to the light, and is the same at
    every pixel (\file{dark\_signal\_characterization.tex}).
  \item \textbf{The stage's light is $n_j P_{ij}$.}  The shape $P$ is
    fixed.  The level $n_j$ changes in a straight line from left to right,
    and so little across one gray chip that it is constant there.  Both
    chips have the same gray, $\rho$.
  \item \textbf{The strip moves little between the photos,} so its shadow
    and its paper move with it and the LEDs see it the same way; lining up
    puts each piece of strip back on its post-image pixel.
  \item \textbf{The dye acts as a filter,} $a = a' t$
    \eqref{eq:dye}.  For dye inside paper this is only roughly true; it
    matters only for what $\iA$ means (step 6), not for the cancellations.
  \item \textbf{A pixel records the light over its area,} so its value is
    the average of $a\,b$ over that area.  The steps below cancel $b$
    exactly where the shadow, or the dye, is even across the pixel.
  \item \textbf{Noise is random}: zero on average, and unrelated between
    the two photos, between pixels, and to the pad's pattern.
  \item \textbf{The gray level $\iG$ carries almost no noise,} because it
    comes from means over the many pixels of each chip.
\end{enumerate}

\subsection*{Running example}

Round numbers, not image data.  Two spots of the pad, fully lit
($b = 1$) and in shadow ($b = 0.75$).  Paper $a' = 0.9$; dye $t = 0.5$, so
$a = 0.45$.  Gray chips $\rho = 0.45$.  Post-image level $n = 1000$,
pre-image $n' = 600$, the same in every column.  Stage shape $P = 1.02$ at
both spots and $1.00$ at the chips.  Dark level $e = e' = 64$\,DN.

%======================================================================
\section{Step 1: remove the dark level (counts)}
\label{sec:dark}

\textbf{In words.}  Take away what the camera reads with no light.

\textbf{Why subtraction, and why first.}  The dark reading is added to the
light (A1), so it is removed by subtraction, in counts.  It must go before
any division: dividing $E$ by anything would scale the dark reading along
with the light, and it could no longer be separated.

\textbf{Algebra.}  Define the light counts
\begin{equation}
  D_{ij} = E_{ij} - e .
  \label{eq:D-def}
\end{equation}
Substituting \eqref{eq:model-pad} for $E_{ij}$ in \eqref{eq:D-def} yields,
for a pad pixel,
\begin{equation}
  D_{ij} = a_{ij}\, b_{ij}\, n_j\, P_{ij} ,
  \label{eq:D}
\end{equation}
and in the same way, from \eqref{eq:model-chip} and \eqref{eq:model-pre},
\begin{equation}
  D_{ij} = \rho\, n_j\, P_{ij} \ \text{(chip)}, \qquad
  D'_{ij} = a'_{ij}\, b_{ij}\, n'_j\, P_{ij}, \qquad
  D'_{ij} = \rho\, n'_j\, P_{ij} \ \text{(chip)}.
  \label{eq:D-others}
\end{equation}

\textbf{Removed:} $e$.  \textbf{Left:} the thought experiment's
$d = a\,b\,n$, with the stage's shape $P$ added.

\textbf{Example.}

\medskip
\begin{center}
\begin{tabular}{@{}lrrrr@{}}
\toprule
 & \multicolumn{2}{c}{post-image} & \multicolumn{2}{c}{pre-image} \\
 & lit & shadow & lit & shadow \\
\midrule
$E$, DN & 523.00 & 408.25 & 614.80 & 477.10 \\
$D = E - e$, DN & 459.00 & 344.25 & 550.80 & 413.10 \\
\bottomrule
\end{tabular}
\end{center}
\noindent For instance $D = 0.45 \times 1 \times 1000 \times 1.02 = 459$.
At the chips, $E = 514$ (post) and $334$ (pre), so $D = 450$ and $270$.

%======================================================================
\section{Step 2: divide by the gray level (counts to reflectance)}
\label{sec:gray}

\textbf{In words.}  The gray chips are read in the same photo as the pad.
Dividing each pad pixel by what a gray chip would read at that pixel
removes the lamp, the exposure time and the shape of the stage's light all
at once.  The result is the pad's \emph{reflectance relative to the gray
chips}: $R = 1$ means ``as bright as a gray chip would be here''.

\textbf{Why here, and why each photo in its own frame.}  The lamp, the
exposure and the stage's shape belong to the stage and to the moment the
photo was taken, not to the strip.  So each photo is divided by its own
chips, at its own pixels, before the pre-image is moved (step 3).

\textbf{Algebra: the level at each chip.}  The chip's level is the mean
of $\iD$ over the chip, divided by the mean of $\iP$ over the same chip:
\begin{equation}
  \gamma_1 = \frac{\avg{g1}{D}}{\avg{g1}{P}} .
  \label{eq:gamma1-def}
\end{equation}
Substituting the chip form of \eqref{eq:D-others} for $D_{ij}$, and taking
$n_j$ as constant across the chip, $n_j = n_1$ (A2), the mean is
\begin{equation}
  \avg{g1}{D} = \rho\, n_1\, \avg{g1}{P} ,
  \label{eq:chipmean}
\end{equation}
and substituting \eqref{eq:chipmean} for $\avg{g1}{D}$ in
\eqref{eq:gamma1-def} yields
\begin{equation}
  \gamma_1 = \rho\, n_1, \qquad \text{and in the same way} \qquad
  \gamma_2 = \rho\, n_2 .
  \label{eq:gamma12}
\end{equation}
The chip's level is the light level there times the chip's gray.  Dividing
by $\avg{g1}{P}$ takes out the stage's shape at the chip.

\textbf{Algebra: the level at each column.}  Between the chips the level
is blended in a straight line.  With $j_1$ and $j_2$ the chips' center
columns,
\begin{equation}
  f_j = \frac{j - j_1}{j_2 - j_1}, \qquad
  \gamma_j = (1 - f_j)\,\gamma_1 + f_j\,\gamma_2 .
  \label{eq:gammaj-def}
\end{equation}
Substituting \eqref{eq:gamma12} for $\gamma_1$ and $\gamma_2$ in
\eqref{eq:gammaj-def} yields
\begin{equation}
  \gamma_j = \rho\,\bigl[(1 - f_j)\, n_1 + f_j\, n_2\bigr] = \rho\, n_j ,
  \label{eq:gammaj}
\end{equation}
because the bracket is the straight line through $n_1$ at column $j_1$ and
$n_2$ at column $j_2$, which by A2 is $n_j$.

\textbf{Algebra: the gray level and the reflectance.}  The gray level
at each pixel is the stage's shape times the column's level, and the
reflectance is the light counts divided by it:
\begin{equation}
  G_{ij} = P_{ij}\, \gamma_j , \qquad
  R_{ij} = \frac{D_{ij}}{G_{ij}} .
  \label{eq:G-R-def}
\end{equation}
Substituting \eqref{eq:gammaj} for $\gamma_j$ in the first of
\eqref{eq:G-R-def} yields
\begin{equation}
  G_{ij} = \rho\, n_j\, P_{ij} ,
  \label{eq:G}
\end{equation}
which is what a gray chip would read at pixel $(i, j)$.  Substituting
\eqref{eq:D} for $D_{ij}$ and \eqref{eq:G} for $G_{ij}$ in the second of
\eqref{eq:G-R-def} yields
\begin{equation}
  R_{ij} = \frac{a_{ij}\, b_{ij}\, n_j\, P_{ij}}{\rho\, n_j\, P_{ij}}
         = \frac{a_{ij}\, b_{ij}}{\rho} .
  \label{eq:R}
\end{equation}
The pre-image goes through the same steps with its own chips, giving
$\gamma'_j = \rho\, n'_j$, $G'_{ij} = \rho\, n'_j P_{ij}$ and, in its own
frame,
\begin{equation}
  R'_{ij} = \frac{a'_{ij}\, b_{ij}}{\rho} .
  \label{eq:Rpre}
\end{equation}

\textbf{Removed:} $n$ (lamp and exposure time) and $P$ (the stage's shape),
inside each photo.  \textbf{Left:} the surface $a$, the shadow $b$, and
the chips' gray $\rho$.\footnote{If the two chips' grays differ, $\rho_1
\ne \rho_2$, then $\gamma_j = (1 - f_j)\rho_1 n_1 + f_j \rho_2 n_2$ is no
longer $\rho\, n_j$, and $R$ is not exactly $a b / \rho$.  The
cancellation in step 4 still holds if the light's left-to-right balance is
the same in both photos, $n'_j = c\, n_j$ for one number $c$: then
$\gamma'_j = c\, \gamma_j$, so $G'_{ij} = c\, G_{ij}$, and in step 4
$R_{ij}/R'_{ij} = (D_{ij}/D'_{ij})(G'_{ij}/G_{ij})
= \bigl(t_{ij}/c\bigr)\, c = t_{ij}$, without knowing $\rho_1$ or
$\rho_2$.}

\textbf{Example.}  $\gamma = 450$ (post) and $270$ (pre) at both chips and
every column, since $P = 1.00$ at the chips.

\medskip
\begin{center}
\begin{tabular}{@{}lrrrr@{}}
\toprule
 & \multicolumn{2}{c}{post-image} & \multicolumn{2}{c}{pre-image} \\
 & lit & shadow & lit & shadow \\
\midrule
$G = P\gamma$, DN & 459.00 & 459.00 & 275.40 & 275.40 \\
$R = D/G$         & 1.000  & 0.750  & 2.000  & 1.500 \\
\bottomrule
\end{tabular}
\end{center}
\noindent Check against \eqref{eq:R}: lit post-image,
$a b/\rho = 0.45 \times 1 / 0.45 = 1.000$; shadowed pre-image,
$0.9 \times 0.75 / 0.45 = 1.500$.  The pad reads about 2 before exposure
because the paper returns about twice what the chips do.

%======================================================================
\section{Step 3: line up the pre-image (an average of light)}
\label{sec:lineup}

\textbf{In words.}  The strip was taken out and put back between the
photos, so a given piece of the pad sits at slightly different pixels in
the two.  Lining up moves the pre-image's reflectance onto the
post-image's grid, so that each post-image pixel is paired with the same
piece of strip.  The result, $\iW$, is the \emph{white}: what this piece
of the pad returned before exposure, one value per pixel.

\textbf{Algebra.}  The piece of strip at post-image pixel $(i, j)$ sat,
in the pre-image, at a position generally between pixels.  Its value is
taken as a weighted average of the four nearest pre-image pixels,
$q = 1 \dots 4$, with weights $\eta_q$ set by how close each one is:
\begin{equation}
  W_{ij} = \sum_{q=1}^{4} \eta_q\, R'_q , \qquad
  \sum_{q=1}^{4} \eta_q = 1 .
  \label{eq:W-def}
\end{equation}

\textbf{Why an arithmetic average of reflectance.}  It stands for the
light a pixel at the in-between position would have collected, and light
adds (rule 1).  $R'$ is light divided by the smooth gray level, so
averaging $R'$ is averaging light.

\textbf{Result.}  With the strip put back in place (A3) and each pixel
recording the light over its area (A5), substituting \eqref{eq:Rpre} into
\eqref{eq:W-def} yields the pre-image's value for the same piece of
strip:
\begin{equation}
  W_{ij} = \frac{a'_{ij}\, b_{ij}}{\rho} .
  \label{eq:W}
\end{equation}

\textbf{Example.}  The example's strip did not move, so $W = R'$: 2.000
(lit) and 1.500 (shadow).

%======================================================================
\section{Step 4: the ratio, without noise}
\label{sec:ratio}

\textbf{In words.}  Divide the post-image by the white.  This is the
thought experiment's ratio \eqref{eq:te-ratio}, now with every factor
accounted for.

\textbf{Algebra.}  Define
\begin{equation}
  T_{ij} = \frac{R_{ij}}{W_{ij}} .
  \label{eq:T-def}
\end{equation}
Substituting \eqref{eq:R} for $R_{ij}$ and \eqref{eq:W} for $W_{ij}$ in
\eqref{eq:T-def} yields
\begin{equation}
  T_{ij} = \frac{a_{ij}\, b_{ij} / \rho}{a'_{ij}\, b_{ij} / \rho}
         = \frac{a_{ij}}{a'_{ij}} ,
  \label{eq:T-a}
\end{equation}
and substituting \eqref{eq:dye} for $a_{ij}$ in \eqref{eq:T-a} yields
\begin{equation}
  T_{ij} = \frac{a'_{ij}\, t_{ij}}{a'_{ij}} = t_{ij} .
  \label{eq:T}
\end{equation}

\textbf{Removed:} $\rho$ (the chips' own gray), $b$ (the shadow), and $a'$
(the paper's own brightness and pattern).  \textbf{Left:} $t$, the
fraction the dye lets through, and nothing else.

\textbf{Where the thought experiment's $n/n'$ went.}  For a strip that did
not move, \eqref{eq:T-def} can be regrouped as
\begin{equation}
  T_{ij} = \frac{D_{ij}/G_{ij}}{D'_{ij}/G'_{ij}}
         = \frac{D_{ij}}{D'_{ij}} \cdot \frac{G'_{ij}}{G_{ij}}
         = \Bigl(t_{ij}\,\frac{n_j}{n'_j}\Bigr) \cdot
           \frac{n'_j}{n_j}
         = t_{ij} ,
  \label{eq:T-check}
\end{equation}
using \eqref{eq:D}, \eqref{eq:D-others}, \eqref{eq:dye} for the first
factor and \eqref{eq:G} for the second.  The first factor is the thought
experiment's ratio \eqref{eq:te-ratio}; the gray chips supply exactly the
$n'/n$ that it left over.

\medskip
\begin{center}
\begin{tabular}{@{}lll@{}}
\toprule
Factor & What it is & Removed in \\
\midrule
$e$ & dark reading & step 1, subtraction \\
$n$ & lamp and exposure time & step 2, division by $\iG$ \\
$P$ & shape of the stage's light & step 2, division by $\iG$ \\
$\rho$ & the chips' own gray & step 4, division by $\iW$ \\
$b$ & shadow & step 4, division by $\iW$ \\
$a'$ & the paper & step 4, division by $\iW$ \\
\midrule
$t$ & the dye & \textbf{kept} \\
\bottomrule
\end{tabular}
\end{center}

\textbf{Example.}  Lit: $T = 1.000 / 2.000 = 0.500$.  Shadow:
$T = 0.750 / 1.500 = 0.500$.  The same stain in and out of the shadow.

%======================================================================
\section{Step 5: noise, and the blended white (method \texorpdfstring{$\MD$}{D})}
\label{sec:noise}

\textbf{In words.}  Every measured value carries a little random noise.
Dividing by $\iW$ removes the shadow and the paper's pattern, but it also
adds the pre-image's noise to the answer.  Smoothing $\iW$ lowers its noise,
but blurs the pattern, which then no longer cancels.  Method $\MD$ uses a
blend of the two, and the algebra below shows that the best blend always
gives some weight to each.

\subsection{The pre-image's noise goes straight into the answer}

Write the measured values as their noise-free values from steps 2 and 3,
$r_{ij} = a_{ij} b_{ij}/\rho$ and $w_{ij} = a'_{ij} b_{ij}/\rho$, times a
small fractional error: $\delta_{ij}$ in the post-image and
$\varepsilon_{ij}$ in the lined-up pre-image,
\begin{equation}
  R_{ij} = r_{ij}\,(1 + \delta_{ij}), \qquad
  W_{ij} = w_{ij}\,(1 + \varepsilon_{ij}) .
  \label{eq:noise}
\end{equation}
($\iG$ adds almost nothing, A7.)  Substituting \eqref{eq:noise} into
\eqref{eq:T-def}, and $r_{ij}/w_{ij} = t_{ij}$ from \eqref{eq:T}, yields
\begin{equation}
  T_{ij} = t_{ij}\, \frac{1 + \delta_{ij}}{1 + \varepsilon_{ij}}
         \approx t_{ij}\,\bigl(1 + \delta_{ij} - \varepsilon_{ij}\bigr),
  \label{eq:T-noise}
\end{equation}
where the approximation keeps only terms of first size (it drops products
of two small errors).  The post-image's noise $\delta$ cannot be avoided.
The pre-image's noise $\varepsilon$ is added to it, in full.

\subsection{Smoothing lowers the noise but leaves a pattern}

Smoothing replaces each value by a weighted average of it and its
neighbors, $m = 0, 1, 2, \dots$ ($m = 0$ is the pixel itself):
\begin{equation}
  \widetilde W_{ij} = \sum_m k_m\, W_m , \qquad \sum_m k_m = 1 .
  \label{eq:Wt-def}
\end{equation}
This is an average of light, so it is done in reflectance (rule 1).
Substituting \eqref{eq:noise} for $W_m$ in \eqref{eq:Wt-def} splits it
into a noise-free part and a noise part:
\begin{equation}
  \widetilde W_{ij} = \sum_m k_m\, w_m \;+\; \sum_m k_m\, w_m\,
  \varepsilon_m .
  \label{eq:Wt-split}
\end{equation}
Define two fractions of $w_{ij}$, one for each part:
\begin{equation}
  \tau_{ij} = \frac{\sum_m k_m\, w_m}{w_{ij}} - 1 , \qquad
  \widetilde\varepsilon_{ij} = \frac{\sum_m k_m\, w_m\,
  \varepsilon_m}{w_{ij}} .
  \label{eq:tau-def}
\end{equation}
$\tau_{ij}$ is how much the smoothing changed the noise-free white at this
pixel: the part of the shadow and paper pattern that smoothing blurred
away.  It is zero where the white is even, and largest where it changes
sharply from pixel to pixel.  $\widetilde\varepsilon_{ij}$ is the noise
left after smoothing: an average of many independent errors (A6), so much
smaller than $\varepsilon_{ij}$.  Substituting \eqref{eq:tau-def} into
\eqref{eq:Wt-split} yields
\begin{equation}
  \widetilde W_{ij} = w_{ij}\,\bigl(1 + \tau_{ij} +
  \widetilde\varepsilon_{ij}\bigr).
  \label{eq:Wt}
\end{equation}
Dividing by $\widetilde W$ instead of $W$ would trade the noise
$\varepsilon$ for the pattern $\tau$.

\subsection{The blend}

Method $\MD$ blends the two, with weights that sum to 1:
\begin{equation}
  \widehat W_{ij} = \lambda\, W_{ij} + (1 - \lambda)\, \widetilde W_{ij},
  \qquad 0 < \lambda < 1 ,
  \label{eq:What-def}
\end{equation}
and the measured fraction the dye let through is now
\begin{equation}
  T_{ij} = \frac{R_{ij}}{\widehat W_{ij}} .
  \label{eq:T5-def}
\end{equation}

\textbf{The blend is itself one average of light.}  Substituting
\eqref{eq:Wt-def} for $\widetilde W_{ij}$ in \eqref{eq:What-def}, and
separating out the pixel itself ($m = 0$), yields
\begin{equation}
  \widehat W_{ij} = \bigl[\lambda + (1 - \lambda)\, k_0\bigr]\, W_{ij}
    + \sum_{m \ne 0} (1 - \lambda)\, k_m\, W_m ,
  \label{eq:What-weights}
\end{equation}
whose weights sum to
$\lambda + (1-\lambda)k_0 + (1-\lambda)(1 - k_0) = \lambda + (1-\lambda) = 1$.
So $\widehat W$ is a single weighted average of the lined-up pre-image:
extra weight on the pixel itself, the rest spread over its neighbors.
Since each $W$ is in turn an average of the pre-image's own pixels
\eqref{eq:W-def}, the whole white is one weighted average of the
pre-image's reflectances.  By rule 1 it is done in reflectance.

\textbf{Algebra.}  Substituting \eqref{eq:noise} for $W_{ij}$ and
\eqref{eq:Wt} for $\widetilde W_{ij}$ in \eqref{eq:What-def} yields
\begin{equation}
  \widehat W_{ij} = w_{ij}\,\bigl[\lambda\,(1 + \varepsilon_{ij})
    + (1 - \lambda)\,(1 + \tau_{ij} + \widetilde\varepsilon_{ij})\bigr]
  = w_{ij}\,\bigl[1 + \lambda\,\varepsilon_{ij}
    + (1 - \lambda)\,(\tau_{ij} + \widetilde\varepsilon_{ij})\bigr],
  \label{eq:What}
\end{equation}
where the 1 comes from $\lambda + (1 - \lambda) = 1$.  Substituting
\eqref{eq:noise} for $R_{ij}$ and \eqref{eq:What} for $\widehat W_{ij}$ in
\eqref{eq:T5-def}, with $r_{ij}/w_{ij} = t_{ij}$, yields the exact result
\begin{equation}
  T_{ij} = t_{ij}\;
    \frac{1 + \delta_{ij}}
         {1 + \lambda\,\varepsilon_{ij}
            + (1 - \lambda)\,(\tau_{ij} + \widetilde\varepsilon_{ij})} .
  \label{eq:T5}
\end{equation}
Without noise or pattern this is $t_{ij}$, as in step 4.  The error terms
are carried into absorbance in step 6, where the best $\lambda$ is found.

\textbf{Example.}  The example has no noise and an even white at each spot,
so $\widehat W = W$ and $T = 0.500$ at both spots, for any $\lambda$.

%======================================================================
\section{Step 6: change to absorbance (once, at the end)}
\label{sec:absorbance}

\textbf{In words.}  Take $-\Lten$ of the fraction the dye let through.
The result is the stain's absorbance: 0 where there is no dye, bigger
where there is more.

\begin{equation}
  A_{ij} = -\Lten T_{ij} .
  \label{eq:A-def}
\end{equation}

\textbf{Why here, and not earlier.}  Every step before this one is either
a division (steps 2 and 4), which gives the same number in either domain
(rule 2), or an average of light (steps 3 and 5), which must be done in
reflectance (rule 1).  So nothing before this point needs absorbance, and
the averages forbid it.  The reasons for absorbance only appear now:
\begin{itemize}[nosep]
  \item \textbf{It measures the amount of dye.}  Light passing through a
    dye keeps a fraction $10^{-\alpha}$, where $\alpha$ grows in
    proportion to the amount of dye (the Beer--Lambert law), so
    $-\Lten t = \alpha$.  For dye inside paper this holds only roughly
    (A4).
  \item \textbf{Two dyes add.}  If yellow lets through $t_u$ and brown
    $t_v$, both together let through $t_u t_v$, and
    $-\Lten (t_u t_v) = -\Lten t_u - \Lten t_v$.  The yellow/brown split
    solves sums, so it needs absorbance.
\end{itemize}
After this point nothing is averaged pixel by pixel, so the domain question
does not come up again.

\textbf{The same number, written as a subtraction.}  Substituting
\eqref{eq:T5-def} for $T_{ij}$ in \eqref{eq:A-def} yields
\begin{equation}
  A_{ij} = -\Lten R_{ij} - \bigl(-\Lten \widehat W_{ij}\bigr) ,
  \label{eq:A-sub}
\end{equation}
the post-image's absorbance minus the white's.  This is rule 2: the final
division and this subtraction are the same step.

\textbf{The result, with its errors.}  Substituting \eqref{eq:T5} for
$T_{ij}$ in \eqref{eq:A-def}, and using $\Lten(xy/z) = \Lten x + \Lten y -
\Lten z$, yields
\begin{equation}
  A_{ij} = -\Lten t_{ij} \;-\; \Lten(1 + \delta_{ij}) \;+\;
    \Lten\bigl[1 + \lambda\,\varepsilon_{ij}
    + (1 - \lambda)(\tau_{ij} + \widetilde\varepsilon_{ij})\bigr] .
  \label{eq:A-exact}
\end{equation}
The first term is $\alpha_{ij}$.  For a small $x$,
$\Lten(1 + x) \approx x / \ln 10 = 0.434\, x$, so to first size
\begin{equation}
  A_{ij} \approx \alpha_{ij} + 0.434\,\bigl[-\delta_{ij}
    + \lambda\,\varepsilon_{ij}
    + (1 - \lambda)(\tau_{ij} + \widetilde\varepsilon_{ij})\bigr] .
  \label{eq:A-first}
\end{equation}
The measured stain is the true amount of dye plus three errors: the
post-image's noise, the pre-image's noise weighted by $\lambda$, and the
blurred pattern (with the small leftover noise) weighted by $1 - \lambda$.

\subsection{The best blend, and why both weights are non-zero}
\label{sec:lambda}

Across the pad, each error has a spread (standard deviation):
$\sigma_\delta$, $\sigma_\varepsilon$, $\sigma_\tau$.  The spread of the
error in \eqref{eq:A-first} follows from the rule that the spreads of
unrelated quantities add as squares (A6: the noises are unrelated to each
other and to the pattern).  Take first the simple case where the smoothing
averages enough pixels that $\widetilde\varepsilon$ can be dropped.  Then
the squared spread of $A - \alpha$, divided by $0.434^2$, is
\begin{equation}
  s^2(\lambda) = \sigma_\delta^2 + \lambda^2\,\sigma_\varepsilon^2
    + (1 - \lambda)^2\,\sigma_\tau^2 .
  \label{eq:spread}
\end{equation}
Its slope with respect to $\lambda$ is
$2\lambda\sigma_\varepsilon^2 - 2(1 - \lambda)\sigma_\tau^2$.  Setting the
slope to zero and solving for $\lambda$ yields the best weight:
\begin{equation}
  \lambda = \frac{\sigma_\tau^2}{\sigma_\varepsilon^2 + \sigma_\tau^2} .
  \label{eq:lambda}
\end{equation}
It is a minimum, because the curvature,
$2(\sigma_\varepsilon^2 + \sigma_\tau^2)$, is positive.  Substituting
\eqref{eq:lambda} for $\lambda$ in \eqref{eq:spread} yields the smallest
spread,
\begin{equation}
  s^2_{\min} = \sigma_\delta^2
    + \frac{\sigma_\varepsilon^2\, \sigma_\tau^2}
           {\sigma_\varepsilon^2 + \sigma_\tau^2} .
  \label{eq:spread-min}
\end{equation}

\noindent What this shows:
\begin{itemize}
  \item \textbf{The weight on the unsmoothed pre-image is the pattern's
    share} of the pre-image's total spread.  A strong pattern calls for
    dividing by the pre-image as it is; strong noise calls for smoothing.
  \item \textbf{Both weights are non-zero} whenever the pre-image has both
    noise ($\sigma_\varepsilon > 0$) and a pattern ($\sigma_\tau > 0$),
    which it always does: \eqref{eq:lambda} is then strictly between 0
    and 1.
  \item \textbf{The blend beats both pure choices.}  At $\lambda = 1$ the
    added spread is $\sigma_\varepsilon^2$; at $\lambda = 0$ it is
    $\sigma_\tau^2$.  The fraction in \eqref{eq:spread-min} is smaller
    than both, because it is $\sigma_\varepsilon^2$ times
    $\sigma_\tau^2/(\sigma_\varepsilon^2 + \sigma_\tau^2) < 1$, and also
    $\sigma_\tau^2$ times
    $\sigma_\varepsilon^2/(\sigma_\varepsilon^2 + \sigma_\tau^2) < 1$.
  \item \textbf{$\lambda$ can be set from measurements:} the pre-image's
    noise $\sigma_\varepsilon$ and the pattern its smoothing removes,
    $\sigma_\tau$.  Where $\sigma_\tau$ differs over the pad, the best
    $\lambda$ differs too.  The algebra holds for any smoothing whose
    weights sum to 1.
\end{itemize}

\noindent\textbf{Example.}  If $\sigma_\varepsilon = 0.010$ and
$\sigma_\tau = 0.006$, then $\lambda = 0.000036/0.000136 = 0.26$, and the
pre-image adds a spread of $0.434\sqrt{0.0000265} = 0.0022$ in absorbance,
against $0.0043$ for $\lambda = 1$ and $0.0026$ for $\lambda = 0$.  (These
spreads are illustrative, not measured.)

\textbf{The exact form.}  Keeping $\widetilde\varepsilon$: if the white
changes little over the smoothing's reach, \eqref{eq:tau-def} gives
$\widetilde\varepsilon_{ij} \approx \sum_m k_m \varepsilon_m$.  Its squared
spread is $h\,\sigma_\varepsilon^2$ with $h = \sum_m k_m^2$, and it shares
the pixel's own error through the weight $k_0$, which adds the term
$2\lambda(1 - \lambda)k_0\sigma_\varepsilon^2$.  Then
\begin{equation}
  s^2(\lambda) = \sigma_\delta^2
    + \sigma_\varepsilon^2\bigl[\lambda^2 + 2\lambda(1-\lambda)k_0
    + (1-\lambda)^2 h\bigr] + (1-\lambda)^2\sigma_\tau^2 ,
  \label{eq:spread-exact}
\end{equation}
and setting its slope to zero yields
\begin{equation}
  \lambda = \frac{\sigma_\tau^2 + \sigma_\varepsilon^2\,(h - k_0)}
                 {\sigma_\tau^2 + \sigma_\varepsilon^2\,(1 - 2k_0 + h)} .
  \label{eq:lambda-exact}
\end{equation}
With $k_0 = h = 0$ this is \eqref{eq:lambda}.  For a bell-shaped smoothing
one pixel wide, $k_0 = 0.159$ and $h = 0.080$; the weight on the
unsmoothed image is then non-zero whenever
$\sigma_\tau^2 > 0.08\,\sigma_\varepsilon^2$, that is, whenever the
pattern's spread is more than 0.28 of the noise's.  (Lining up shares
noise between neighboring pixels, which changes $h$ and $k_0$ slightly but
not the form of the result.)  Both formulas were checked numerically by
finding the minimum of \eqref{eq:spread} and \eqref{eq:spread-exact} on a
fine grid of $\lambda$.

\subsection{All channels}

Steps 1--6 are done for each channel separately.  Putting each channel's
$\iA$ side by side, one row per pad pixel and one column per channel, gives
the table $\tblA$ that the yellow/brown split works on.

%======================================================================
\section{The chain on one page}
\label{sec:summary}

\begin{align*}
  &\text{1. counts:}       & D_{ij} &= E_{ij} - e
     &&= a_{ij}\, b_{ij}\, n_j\, P_{ij} \\
  &\text{2. reflectance:}  & R_{ij} &= D_{ij} \big/ (P_{ij}\gamma_j)
     &&= a_{ij}\, b_{ij} / \rho \\
  &\text{2. pre-image:}    & R'_{ij} &= D'_{ij} \big/ (P_{ij}\gamma'_j)
     &&= a'_{ij}\, b_{ij} / \rho \quad\text{(own frame)} \\
  &\text{3. line up:}      & W_{ij} &= \textstyle\sum_q \eta_q R'_q
     &&= a'_{ij}\, b_{ij} / \rho \\
  &\text{5. blend:}        & \widehat W_{ij} &= \lambda W_{ij}
     + (1-\lambda)\textstyle\sum_m k_m W_m
     &&= w_{ij}\,[1 + \text{small errors, \eqref{eq:What}}] \\
  &\text{4, 5. ratio:}     & T_{ij} &= R_{ij} / \widehat W_{ij}
     &&\approx t_{ij} \\
  &\text{6. absorbance:}   & A_{ij} &= -\Lten T_{ij}
     &&\approx \alpha_{ij}
\end{align*}

\medskip
\begin{tabular}{@{}lL{8.4cm}@{}}
\toprule
Assumption & Used in \\
\midrule
A1 dark adds, even & step 1 \\
A2 stage light $n_j P_{ij}$, one gray & step 2 \\
A3 strip moves little & step 3 (and $b$ the same in both photos) \\
A4 dye is a filter & step 4 ($a = a't$); step 6 (what $\iA$ means) \\
A5 pixel averages its area & step 3, step 4 \\
A6 noise random, unrelated & step 5, the best $\lambda$ \\
A7 $\iG$ nearly noise-free & step 5 \\
\bottomrule
\end{tabular}

\medskip
\noindent There is no single white number anywhere: every pixel is divided
by its own white, $\widehat W_{ij}$, taken from the same strip before
exposure.

%======================================================================
\section{Where this differs from the code today}
\label{sec:code}

\begin{itemize}[nosep]
  \item The smoothing and the blend are averages of light, so here they are
    done in reflectance.  The study code (\file{shading\_review.py}) does
    them in absorbance.  The two agree to first size and differ by small
    second-size terms; the difference is largest where the white changes
    sharply.
  \item Absorbance is taken once, at the end, instead of before the
    shading removal.
  \item There is no scalar white (\file{R\_white}).
\end{itemize}

\appendix
\section{Symbols in \texttt{analysis\_chain.tex}}

\begin{tabular}{@{}lll@{}}
\toprule
This document & \file{analysis\_chain.tex} & Meaning \\
\midrule
$\iD$ & $\mathbf{F}$ & light counts \\
$\iW$ & $\mathbf{S}$ & pre-image, lined up \\
$\hW$ & $\widehat{\mathbf{S}}$ & estimated white \\
$\iT$ & $\mathbf{T}$ & corrected reflectance \\
$\iA$ & $\mathbf{C}$ & stain absorbance \\
--- & $\mathbf{A}$, $\mathbf{B}$, $\widetilde{\mathbf{B}}$,
  $\widehat{\mathbf{B}}$ & not needed: nothing is averaged in absorbance \\
\bottomrule
\end{tabular}

\end{document}
